%% ============================================================
%% Episode 05 - Farsi Podcast Script (Translation)
%% Source: specialized_advisor_report_complete/markdown/05_lyapunov_stability.md
%% Series: Advisor Progress Report - Deep Dive
%% Compiler: XeLaTeX ONLY
%% Font: B Nazanin
%% CRITICAL: xepersian must always be loaded LAST
%% ============================================================

\documentclass[12pt, a4paper]{article}

%% ============================================================
%% COMPATIBILITY SHIMS
%% ============================================================
\makeatletter
\providecommand\IfClassLoadedT[2]{\@ifclassloaded{#1}{#2}{}}
\providecommand\IfClassLoadedF[2]{\@ifclassloaded{#1}{}{#2}}
\providecommand\IfPackageLoadedT[2]{\@ifpackageloaded{#1}{#2}{}}
\providecommand\IfPackageLoadedF[2]{\@ifpackageloaded{#1}{}{#2}}
\providecommand\IfFileLoadedTF[3]{\@ifpackageloaded{#1}{#2}{#3}}
\providecommand\IfFileLoadedT[2]{\@ifpackageloaded{#1}{#2}{}}
\providecommand\IfFileLoadedF[2]{\@ifpackageloaded{#1}{}{#2}}
\makeatother
\usepackage{expl3}
\ExplSyntaxOn
\cs_if_exist:NF \prop_gput_if_not_in:Nnn
  {
    \cs_new:Npn \prop_gput_if_not_in:Nnn #1#2#3
      { \prop_if_in:NnF #1 {#2} { \prop_gput:Nnn #1 {#2} {#3} } }
  }
\ExplSyntaxOff
%% ============================================================

\usepackage[
  a4paper,
  top=2.5cm, bottom=2.5cm,
  left=2.8cm, right=2.8cm,
  headheight=25pt
]{geometry}

\usepackage{xcolor}
\definecolor{seriesblue}{RGB}{25, 80, 150}
\definecolor{audionote}{RGB}{130, 50, 15}
\definecolor{sectiongray}{RGB}{75, 75, 75}
\definecolor{audiotan}{RGB}{252, 243, 233}

\usepackage{amsmath}
\usepackage{amssymb}
\usepackage{enumitem}
\usepackage{setspace}
\setstretch{1.8}
\usepackage{mdframed}
\usepackage{titlesec}
\usepackage{booktabs}

%% ============================================================
%% xepersian MUST be loaded LAST
%% ============================================================
\usepackage{xepersian}
\settextfont[Scale=1.05]{B Nazanin}
\setlatintextfont{Times New Roman}

\newcommand{\dash}{\lr{---}}

%% ============================================================
%% Page style -- savebox approach
%% ============================================================
\newsavebox{\hdrLeft}
\newsavebox{\hdrRight}
\AtBeginDocument{%
  \savebox{\hdrLeft}{\normalfont\small\color{sectiongray}%
    گزارش پيشرفت مشاور / بررسی عمیق}%
  \savebox{\hdrRight}{\normalfont\small\color{sectiongray}%
    قسمت پنجم: پايداری لياپانوف}%
}
\makeatletter
\def\ps@podcast{%
  \def\@oddhead{\usebox{\hdrLeft}\hfill\usebox{\hdrRight}}%
  \let\@evenhead\@oddhead
  \def\@oddfoot{\normalfont\hfil\thepage\hfil\relax}%
  \let\@evenfoot\@oddfoot
  \let\@mkboth\markboth
}
\makeatother
\pagestyle{podcast}

%% ============================================================
%% Section formatting
%% ============================================================
\titleformat{\section}
  {\large\bfseries\color{seriesblue}}
  {}
  {0pt}
  {}
  [\vspace{-2pt}{\color{seriesblue}\rule{\linewidth}{0.7pt}}\vspace{2pt}]

\titleformat{\subsection}
  {\normalsize\bfseries\color{sectiongray}}
  {}
  {0pt}
  {}

\setcounter{secnumdepth}{0}

\setlist[itemize]{
  itemsep=4pt,
  topsep=5pt,
  parsep=0pt,
  label=$\bullet$
}

%% ============================================================
\begin{document}
%% ============================================================

\thispagestyle{empty}

\vspace*{1.2cm}

{\centering
{\color{seriesblue}\fontsize{24}{30}\selectfont\bfseries قسمت پنجم}

\vspace{0.5cm}
{\fontsize{15}{20}\selectfont\bfseries پايداری لياپانوف: اثبات رياضی کارکرد سيستم}

\vspace{0.9cm}
{\color{seriesblue}\rule{0.55\linewidth}{1pt}}
\vspace{0.7cm}

{\normalsize
\begin{tabular}{rl}
  \textbf{مجموعه:}     & گزارش پيشرفت مشاور \dash{} بررسی عمیق \\[5pt]
  \textbf{مدت:}        & ۱۰ تا ۱۲ دقیقه \\[5pt]
  \textbf{راوی:}       & مجری تنها \\[5pt]
  \textbf{منابع:}      & بخش‌های ۴.۱ تا ۴.۷ گزارش مشاور \\
\end{tabular}
}

\vspace{0.9cm}
{\color{seriesblue}\rule{0.55\linewidth}{0.4pt}}
\par}

\vspace{1.2cm}

%% ============================================================
\section{افتتاحیه: چرا لياپانوف؟}

نتايج شبيه‌سازی که نشان می‌دهند يک کنترل‌گر کار می‌کند يک اثبات نيستند. آن‌ها شواهد هستند. يک مشاور خواهد پرسيد: آيا می‌توانيد رياضياتاً ثابت کنيد کنترل‌گر سيستم را پايدار می‌کند \dash{} نه فقط برای شرايط اوليه خاصی که آزمايش کرديد، بلکه به طور کلی؟

نظريه پايداری لياپانوف دقيقاً همين را فراهم می‌کند. ايده اصلی زيباست: يک تابع \lr{V} از حالت پيدا کنيد که مثل انرژی رفتار کند. بايد مثبت باشد \dash{} مثل انرژی، هميشه غيرمنفی است. و مشتق زمانی‌اش در طول مسيرهای سيستم بايد منفی باشد \dash{} مثل انرژی در حال اتلاف از يک سيستم ميراشده. اگر بتوانيد چنين تابعی پيدا کنيد و نشان دهيد مشتق منفی است، پايداری بدون نياز به حل معادلات ديفرانسيل ثابت می‌شود.

اثبات‌های کامل در مستندات نظری هستند \dash{} ۱۴۲۷ خط، نه قضيه، پنج لم. اين قسمت نتايج کليدی گزارش را پوشش می‌دهد.

%% ============================================================
\section{توضيح مهم: بتا و اتا چيستند؟}

قبل از اثبات‌های کنترل‌گر، گزارش يک نکته را بيان می‌کند که ارزش دارد به صراحت گفته شود.

نابرابری‌های پايداری شامل ثابت‌هايی مثل بتا، اتا، سی-يک، سی-دو، آلفا-يک، آلفا-دو هستند. اينها اعداد ثابت نيستند. آن‌ها کميت‌هايی هستند که به حالت يا به بهره بستگی دارند.

بتا مساوی است با \lr{L} ضربدر معکوس \lr{M} ضربدر \lr{B} \dash{} همان حاصل‌ضربی که در بررسی درجه نسبی ديديم. اين به پيکربندی آونگ فعلی از طريق معکوس \lr{M} بستگی دارد. با نوسان آونگ‌ها، بتا تغيير می‌کند. اثبات‌ها نشان می‌دهند بتا در طول همه مسيرهای فيزيکی معتبر مثبت می‌ماند \dash{} که می‌ماند، چون \lr{M} هميشه مثبت معين است.

اتا مساوی است با \lr{K} منهای \lr{d-bar}، که \lr{K} بهره سوئيچينگ و \lr{d-bar} کران اغتشاش است. اين با طراحی مثبت است: \lr{K} را ملزم می‌کنيم از \lr{d-bar} بيشتر باشد.

درک اين موضوع برای سوالات مشاور مهم است. «مقدار بتا چقدر است؟» سوال درستی نيست. پاسخ درست: بتا يک اسکالر وابسته به پيکربندی است که در طول حرکت مثبت می‌ماند، و اثبات اين را به صراحت تأييد می‌کند.

%% ============================================================
\section{کنترل مُد لغزشی کلاسيک: پايداری نمايی مجانبی}

تابع لياپانوف: \lr{V} مساوی است با نصف سيگما به توان دو.

اين ساده‌ترين کانديدای ممکن است \dash{} مجذور مقدار سطح لغزشی. وقتی سيگما غيرصفر است مثبت است، و دقيقاً روی سطح صفر است.

مشتق زمانی: \lr{V-dot} مساوی سيگما ضربدر سيگما-دات. با جايگزين کردن ديناميک‌های حلقه بسته، گزارش نشان می‌دهد:

\lr{V-dot} کمتر يا مساوی است با منفی بتا ضربدر اتا ضربدر قدرمطلق سيگما، منهای بتا ضربدر کا-دی ضربدر سيگما به توان دو.

هر دو جمله خارج از لايه مرزی منفی هستند. جمله اول \dash{} خطی در قدرمطلق سيگما \dash{} دسترسی را هدايت می‌کند. جمله دوم \dash{} درجه دو در سيگما \dash{} همگرايی نمايی را فراهم می‌کند.

نوع پايداری: مجانبی نمايی. تصديق عددی: ۹۶.۲ درصد نمونه‌های شبيه‌سازی خارج از لايه مرزی \lr{V-dot} کمتر از صفر دارند. ۳.۸ درصد باقيمانده در طول فاز دسترسی اوليه رخ می‌دهند که در آن مسير از لايه مرزی عبور می‌کند. اين رفتار مورد انتظار است.

%% ============================================================
\section{الگوريتم ابر-پيچشی: همگرايی زمان محدود}

تابع لياپانوف: \lr{V-STA} مساوی است با قدرمطلق سيگما، به علاوه يک دوم \lr{K-two} ضربدر \lr{z} به توان دو، که \lr{z} مساوی \lr{u-two} مؤلفه انتگرالی است.

اين ساختار لياپانوف مورنو-اوسوريو برای مُدهای لغزشی مرتبه دوم است.

نتيجه کليدی: وقتی \lr{L} \dash{} ثابت ليپشيتز مشتق اغتشاش \dash{} صفر باشد، \lr{V-dot-STA} کاملاً منفی است و همگرايی زمان محدود تضمين می‌شود. برخلاف کنترل کلاسيک که به صورت مجانبی همگرا می‌شود، \lr{STA} سيگما را در زمان محدود \lr{T-f} دقيقاً به صفر می‌راند.

شرايط بهره مورد نياز: \lr{K-one} بايد از دو ضربدر ريشه دو ضربدر \lr{d-bar} تقسيم بر بتا بيشتر باشد، و \lr{K-two} بايد از \lr{d-bar} تقسيم بر بتا بيشتر باشد، و \lr{K-one} بايد از \lr{K-two} بيشتر باشد.

بهره‌های اسمی \lr{K-one} مساوی ۸ و \lr{K-two} مساوی ۴ اين شرايط را ارضا می‌کنند. بهره‌های استحکام \lr{K-one} مساوی ۲.۰۲ و \lr{K-two} مساوی ۶.۶۷ شرط \lr{K-one} بزرگ‌تر از \lr{K-two} را ارضا نمی‌کنند \dash{} که مسئله باز مستند است.

نتيجه تجربی: همگرايی در ۱.۸۲ ثانيه برای شرايط اوليه معمول، ۱۶ درصد سريع‌تر از کنترل کلاسيک در ۲.۱۵ ثانيه.

%% ============================================================
\section{کنترل مُد لغزشی تطبيقی: مجانبی با بهره کران‌دار}

تابع لياپانوف: \lr{V-ad} مساوی است با نصف سيگما به توان دو، به علاوه يک دوم گاما ضربدر \lr{K-tilde} به توان دو، که \lr{K-tilde} اختلاف بين بهره تطبيقی فعلی و بهره ايده‌آل \lr{K-star} است.

اين يک تابع لياپانوف مرکب است \dash{} شامل هم خطای حالت از طريق سيگما و هم خطای تخمين پارامتر از طريق \lr{K-tilde}. اين رويکرد استاندارد برای اثبات‌های پايداری کنترل تطبيقی است.

نوع پايداری: همگرايی مجانبی حالت با بهره تطبيقی کران‌دار. تصديق عددی: ۱۰۰ درصد نمونه‌های شبيه‌سازی \lr{K(t)} را در محدوده ۰.۱ تا ۱۰۰ نشان می‌دهند.

%% ============================================================
\section{هيبريد تطبيقی ابر-پيچشی: پايداری ورودی به حالت}

تابع لياپانوف: \lr{V-hyb} مساوی است با نصف سيگما به توان دو، به علاوه کا-يک-تيلدا به توان دو تقسيم بر دو گاما-يک، به علاوه کا-دو-تيلدا به توان دو تقسيم بر دو گاما-دو، به علاوه نصف \lr{u-int} به توان دو.

اين ساختار مرکب را برای شامل کردن هم خطاهای بهره تطبيقی و هم حالت انتگرالی گسترش می‌دهد.

مشتق زمانی: \lr{V-dot-hyb} کمتر يا مساوی است با منفی آلفا-يک ضربدر \lr{V-hyb}، به علاوه آلفا-دو ضربدر نُرم بردار اغتشاش \lr{w}. اين پايداری ورودی به حالت يا \lr{ISS} است. سيستم به صورت غيرمشروط پايدار نيست \dash{} جمله آلفا-دو ضربدر نُرم اغتشاش می‌تواند مثبت باشد. اما اگر اغتشاش کران‌دار باشد، حالت کران‌دار می‌ماند. \lr{ISS} يک مفهوم پايداری ضعيف‌تر اما همچنان دقيق است که برای سيستم‌هايی مناسب است که در معرض اغتشاشات مداوم هستند.

%% ============================================================

\begin{mdframed}[
  backgroundcolor=audiotan,
  linecolor=audionote!70,
  linewidth=1pt,
  innerleftmargin=12pt,
  innerrightmargin=12pt,
  innertopmargin=8pt,
  innerbottommargin=8pt
]
\textbf{\color{audionote}[يادداشت صوتی:]}
جدول خلاصه: کنترل کلاسيک \dash{} تابع نصف سيگما به توان دو \dash{} مجانبی نمايی \dash{} ۲.۱۵ ثانيه. الگوريتم ابر-پيچشی \dash{} قدرمطلق سيگما به علاوه \lr{z} به توان دو \dash{} زمان محدود \dash{} ۱.۸۲ ثانيه. تطبيقی \dash{} مرکب با بهره \dash{} مجانبی کران‌دار \dash{} ۲.۳۵ ثانيه. هيبريد تطبيقی \dash{} مرکب گسترده \dash{} \lr{ISS} \dash{} ۱.۹۵ ثانيه. ترتيب همگرايی از سريع‌ترين تا کندترين: \lr{STA}، هيبريد، کلاسيک، تطبيقی.
\end{mdframed}

\vspace{0.6cm}

%% ============================================================
\section{نتيجه‌گيری}

آناليز لياپانوف استخوان‌بندی نظری پايان‌نامه است. هر چهار کنترل‌گر اثبات‌های کامل دارند: کنترل کلاسيک با پايداری نمايی مجانبی، \lr{STA} با همگرايی زمان محدود قوی‌ترين تضمين، تطبيقی با بهره کران‌دار، و هيبريد با \lr{ISS} مناسب برای اغتشاشات مداوم.

در قسمت بعدی: از نظريه به آزمايش می‌رويم \dash{} چگونه معيارهای مونت‌کارلو طراحی و اجرا شدند؟

%% ============================================================
\vspace{1.5cm}
\begin{center}
{\color{seriesblue}\rule{0.45\linewidth}{0.4pt}}\\[8pt]
{\small\color{sectiongray}
منابع گزارش: بخش‌های ۴.۱ تا ۴.۷ \dash{} معادلات \lr{eq:lyap\_classical, eq:lyap\_classical\_dot, eq:lyap\_sta, eq:sta\_conditions}
}
\end{center}

\end{document}
